\documentclass[11pt,a4paper]{article}
% Drake_preamble.tex - LaTeX Preamble Template
% 作者: 謝慕揚, MD, PhD, FESC
% 
% 使用說明:
% 1. 表格請使用 longtable，不要有垂直分隔線 |
% 2. 表格內換行使用 \newline，不要使用 \\
% 3. 藥物名稱、檢驗、檢查請維持英文
% 4. Reference 格式請使用 \href 加上驗證過的 DOI

% Including essential packages for Chinese typesetting
\usepackage{xeCJK}
\usepackage{fontspec}
% Configuring Chinese font with Noto Serif CJK TC, fallback to SimSun
\IfFontExistsTF{Noto Serif CJK TC}{
  \setCJKmainfont{Noto Serif CJK TC}
  \setCJKsansfont{Noto Serif CJK TC}
  \setCJKmonofont{Noto Serif CJK TC}
}{\setCJKmainfont{SimSun}
  \setCJKsansfont{SimSun}
  \setCJKmonofont{SimSun}
}

% Including other necessary packages
\usepackage{geometry}
\usepackage{titlesec}
\usepackage{enumitem}
\usepackage{xcolor}
\usepackage{colortbl}
\usepackage{tcolorbox}
\usepackage{booktabs}
\usepackage{longtable}
\usepackage{array}
\usepackage{graphicx}
\usepackage{tikz}
\usepackage{fancyhdr}
\usepackage{multicol}
\usepackage{datetime2}
\usepackage{hyperref}
\usepackage{mdframed}
\usepackage{amsmath}
\usepackage{amssymb}
\usepackage{multirow}

\hypersetup{
    colorlinks=true,
    linkcolor=emergency_blue,
    citecolor=emergency_blue,
    urlcolor=emergency_blue,
    pdfborder={0 0 0}
}

% Defining page geometry
\geometry{
  top=2.5cm,
  bottom=2.5cm,
  left=2cm,
  right=2cm,
  headheight=25pt
}

% Adjusting topmargin to compensate for increased headheight
\addtolength{\topmargin}{-10pt}

% Defining colors
\definecolor{emergency_red}{RGB}{186,24,27}
\definecolor{emergency_blue}{RGB}{0,114,188}
\definecolor{emergency_gray}{RGB}{120,120,120}
\definecolor{highlight_yellow}{RGB}{255,250,205}

% Setting title formats
\titleformat{\section}[block]{\Large\bfseries\color{emergency_red}}{\thesection}{10pt}{}
\titleformat{\subsection}[block]{\large\bfseries\color{emergency_blue}}{\thesubsection}{8pt}{}
\titleformat{\subsubsection}[block]{\normalsize\bfseries\color{emergency_gray}}{\thesubsubsection}{6pt}{}

% Defining custom environment for important notes
\newmdenv[
    linecolor=emergency_red,
    backgroundcolor=highlight_yellow,
    linewidth=2pt,
    topline=true,
    bottomline=true,
    leftline=true,
    rightline=true,
    innertopmargin=10pt,
    innerbottommargin=10pt,
    innerleftmargin=10pt,
    innerrightmargin=10pt
]{importantbox}

% Defining note command with tcolorbox
\newcommand{\note}[1]{
    \par
    \noindent
    \begin{tcolorbox}[
        colback=emergency_blue!5!white,
        colframe=emergency_blue,
        title=\textbf{重點提醒},
        fonttitle=\bfseries,
        boxrule=0.5pt,
        arc=3pt,
        left=6pt,
        right=6pt,
        top=6pt,
        bottom=6pt
    ]
    #1
    \end{tcolorbox}
}

% Key findings box
\newcommand{\keyfinding}[1]{
    \par
    \noindent
    \begin{tcolorbox}[
        colback=yellow!10!white,
        colframe=orange!75!black,
        title=\textbf{核心發現摘要},
        fonttitle=\bfseries,
        boxrule=0.5pt,
        arc=3pt
    ]
    #1
    \end{tcolorbox}
}

% Defining custom date format for ISO 8601
\DTMnewdatestyle{iso}{
  \renewcommand{\DTMdisplaydate}[4]{\number##1-\number##2-\number##3}
  \renewcommand{\DTMDisplaydate}[4]{\number##1-\number##2-\number##3}
}
\DTMsetdatestyle{iso}
\newcommand{\mydate}{\DTMtoday}


% Custom header
\pagestyle{fancy}
\fancyhf{}
\fancyhead[L]{\textcolor{emergency_red}{CHEST 教學講義}}
\fancyhead[R]{\textcolor{emergency_blue}{ILD-PH 診斷與治療}}
\fancyfoot[C]{\textcolor{emergency_gray}{\thepage}}
\renewcommand{\headrulewidth}{1pt}
\renewcommand{\headrule}{\hbox to\headwidth{\color{emergency_red}\leaders\hrule height \headrulewidth\hfill}}

\begin{document}

%============================================================
% Title Page
%============================================================
\begin{center}
{\LARGE\bfseries\textcolor{emergency_red}{CHEST Journal}}\\[0.5cm]
{\Large\textbf{Pulmonary Hypertension Associated With Interstitial Lung Diseases}}\\[0.3cm]
{\large 間質性肺病相關肺高壓：診斷與治療實務}\\[0.5cm]
{\normalsize \href{https://doi.org/10.1016/j.chest.2025.07.4107}{CHEST 2026;169(1):220-229. DOI: 10.1016/j.chest.2025.07.4107}}\\[0.3cm]
{\normalsize 整理：謝慕揚 MD, PhD, FESC \quad 日期：\mydate}
\end{center}

\vspace{0.5cm}

%============================================================
\section{前言}
%============================================================

間質性肺病 (interstitial lung disease, ILD) 涵蓋多種以肺實質發炎與纖維化為特徵的肺部疾病。肺高壓 (pulmonary hypertension, PH) 常見於纖維性 ILD 患者，診斷與治療皆具挑戰性。

\keyfinding{
\begin{itemize}
    \item ILD-PH 患者死亡率較單純 PAH 高 3 倍以上
    \item 嚴重 PH 約發生於 5-10\% 的纖維性 ILD 患者
    \item Combined pulmonary fibrosis with emphysema (CPFE) 患者 PH 發生率可達 50\%
    \item Inhaled treprostinil (iTRE) 是目前唯一有 RCT 證據支持的 ILD-PH 治療藥物
\end{itemize}
}

%============================================================
\section{ILD-PH 定義與流行病學}
%============================================================

\subsection{血行動力學定義}

根據 2022 ESC/ERS 指引，precapillary PH 定義為：
\begin{itemize}
    \item Mean pulmonary artery pressure (mPAP) $>$ 20 mmHg
    \item Pulmonary capillary wedge pressure (PCWP) $\leq$ 15 mmHg
    \item Pulmonary vascular resistance (PVR) $>$ 2 Wood units (WU)
\end{itemize}

\subsection{流行病學}

\begin{longtable}{p{5cm}p{9cm}}
\toprule
\textbf{ILD 類型} & \textbf{PH 發生率與特點} \\
\midrule
\endfirsthead
\toprule
\textbf{ILD 類型} & \textbf{PH 發生率與特點} \\
\midrule
\endhead
Idiopathic pulmonary fibrosis (IPF) & 嚴重 PH 約 5-10\%\newline 預後不良 \\
\midrule
CPFE & PH 發生率可達 50\%\newline 嚴重 PH 比例較其他 ILD 高 \\
\midrule
Connective tissue disease-associated ILD & 需排除 PAH (Group 1)\newline 可能有多重機轉 \\
\bottomrule
\end{longtable}

%============================================================
\section{診斷流程}
%============================================================

\subsection{臨床懷疑 ILD-PH 的線索}

\begin{longtable}{p{4cm}p{10cm}}
\toprule
\textbf{檢查項目} & \textbf{提示 PH 的發現} \\
\midrule
\endfirsthead
\toprule
\textbf{檢查項目} & \textbf{提示 PH 的發現} \\
\midrule
\endhead
Transthoracic echocardiography (TTE) & Right atrial area $>$ 18 cm$^2$\newline RV-LV end-diastolic ratio $>$ 1\newline LV eccentricity index $>$ 1.1\newline TR jet velocity $>$ 2.8 m/s\newline TAPSE $<$ 1.8 cm \\
\midrule
Pulmonary function test (PFT) & Isolated decrease in DLCO\newline Elevated FVC/DLCO ratio \\
\midrule
Chest CT & PA diameter $>$ 2.9 cm\newline PA:Aorta ratio $>$ 1\newline RV dilation or hypertrophy \\
\midrule
6-minute walk test (6MWT) & Drop in distance\newline O$_2$ saturation nadir\newline Abnormal heart rate recovery \\
\midrule
Biomarkers & Elevated BNP or NT-proBNP \\
\bottomrule
\end{longtable}

\note{
TTE 在 ILD 患者因肺部過度充氣導致 acoustic window 不佳，RVSP 在高達 50\% 患者無法測量。即使可測量，與 RHC 測得的 systolic PA pressure 相關性也不佳。
}

\subsection{Right Heart Catheterization (RHC)}

RHC 是確診 ILD-PH 的必要檢查：

\begin{itemize}
    \item 應在患者臨床穩定時進行（非 ILD 急性發作期）
    \item 約 20\% ILD-PH 患者有 post-capillary 或 combined pre- and post-capillary PH
    \item 有 left heart disease 風險因子者，考慮 provocative testing (fluid or exercise challenge)
    \item PCWP 測量困難時，可考慮 left heart catheterization 測量 LVEDP
    \item \textbf{Vasoreactivity testing 不建議用於 ILD-PH}
\end{itemize}

%============================================================
\section{鑑別診斷}
%============================================================

ILD 患者發生 PH 需排除其他原因：

\begin{itemize}
    \item \textbf{Group 2 PH}：Left heart disease（共同危險因子如抽菸、年齡）
    \item \textbf{Group 4 PH}：Chronic thromboembolic PH（ILD 患者 VTE 風險增加）
    \item \textbf{Group 1 PAH}：Connective tissue disease、HIV infection
\end{itemize}

%============================================================
\section{治療策略}
%============================================================

\subsection{基礎治療（所有 ILD-PH 患者）}

\begin{itemize}
    \item 戒菸
    \item 補充氧氣（若 rest、activity 或 sleep 時低血氧）
    \item ILD 特異性治療（antifibrotic therapy、immunosuppressive therapy）
    \item Inhalers（若合併 obstructive lung disease）
    \item 睡眠呼吸障礙及慢性高碳酸血症之非侵襲性呼吸器支持
    \item 肺復健
    \item 疫苗接種
    \item 利尿劑（若有體液過載）
    \item \textbf{肺移植評估}
\end{itemize}

\subsection{PH 特異性藥物治療}

\begin{longtable}{p{4cm}p{10cm}}
\toprule
\textbf{藥物類別} & \textbf{ILD-PH 使用建議} \\
\midrule
\endfirsthead
\toprule
\textbf{藥物類別} & \textbf{ILD-PH 使用建議} \\
\midrule
\endhead
\textbf{Inhaled treprostinil (iTRE)} & \textcolor{emergency_blue}{唯一有 RCT 證據}\newline INCREASE trial：改善 6MWT、BNP、減少 clinical worsening\newline 最佳效果：$\geq$ 9 breaths QID、baseline PVR $\geq$ 4 WU\newline \textbf{注意}：CPFE 以 emphysema 為主者應避免使用 \\
\midrule
Soluble guanylate cyclase stimulator (riociguat) & \textcolor{emergency_red}{禁忌於 IIP-PH}\newline RISE-IIP trial 因嚴重不良事件及死亡率增加而提前終止 \\
\midrule
Endothelin receptor antagonist (ambrisentan) & \textcolor{emergency_red}{禁忌於 IPF-PH}\newline ARTEMIS trial 因呼吸住院及疾病惡化風險增加而提前終止 \\
\midrule
PDE-5 inhibitor (sildenafil, tadalafil) & 無 RCT 證據\newline Registry 及 retrospective data 顯示可能改善血行動力學及存活\newline 當 iTRE 不可得時，可在專家中心個別化使用 \\
\midrule
Parenteral prostacyclin & 僅用於嚴重 RV failure\newline 需在專家中心進行充分風險效益討論 \\
\bottomrule
\end{longtable}

\subsection{INCREASE Trial 重點}

\begin{itemize}
    \item \textbf{設計}：雙盲 RCT，ILD-PH 患者 1:1 分配至 iTRE 或 placebo
    \item \textbf{受試者}：約 50\% IPF、25\% CTD-ILD、25\% CPFE
    \item \textbf{主要結果}：
    \begin{itemize}
        \item 6MWT distance 改善
        \item BNP 下降
        \item Clinical worsening 減少
    \end{itemize}
    \item \textbf{Post-hoc 分析}：
    \begin{itemize}
        \item 穩定 FVC
        \item 減少 ILD exacerbation 風險
    \end{itemize}
    \item \textbf{Subgroup 分析}：CPFE-PH 患者 6MWT 獲益不如其他 ILD 類型
\end{itemize}

\note{
PERFECT trial（iTRE 用於 COPD-PH）因嚴重不良事件及死亡率增加而提前終止，特別是 DLCO $<$ 25\% predicted 者。因此建議 iTRE 保留給以纖維化為主、emphysema 負擔較輕的 CPFE 患者。
}

%============================================================
\section{「輕度 ILD」合併顯著 PH}
%============================================================

部分患者有相對輕微的 ILD 但顯著的 PH，特徵如下：

\begin{itemize}
    \item 年齡較大、男性居多
    \item DLCO 較低
    \item 抽菸史較多
    \item 預後類似 ILD-PH，但治療反應不如 IPAH
    \item 建議轉介至 PH center 管理
\end{itemize}

「輕度 ILD」定義尚未標準化（建議：CT 上纖維化 $<$ 20\% 或 FVC $>$ 70\%）。

%============================================================
\section{Key Points}
%============================================================

\begin{enumerate}
    \item 嚴謹的 phenotyping（PFT、HRCT）對確保正確診斷與治療至關重要
    \item 評估應在患者臨床穩定時進行，因 ILD 急性發作可暫時升高肺動脈壓
    \item Post-capillary PH 高風險者，RHC 時考慮 provocative maneuvers
    \item 治療為多面向：基礎 ILD 治療、矯正低血氧、處理共病、考慮肺移植、決定 PH 治療
    \item \textbf{iTRE 對纖維性 ILD-PH 有益}，但應避免用於 emphysema 為主的 CPFE-PH
    \item PDE-5i 及 parenteral prostacyclin 可在專家中心個別化使用
    \item Soluble guanylate cyclase stimulator 及 ERA 在某些 ILD 類型禁忌
\end{enumerate}

%============================================================
\section{參考文獻}
%============================================================

\begin{enumerate}
    \item Humbert M, Kovacs G, Hoeper MM, et al. 2022 ESC/ERS guidelines for the diagnosis and treatment of pulmonary hypertension. \href{https://doi.org/10.1093/eurheartj/ehac237}{\textit{Eur Heart J}. 2022;43(38):3618-3731.}

    \item Shlobin OA, Adir Y, Barbera JA, et al. Pulmonary hypertension associated with lung diseases. \href{https://doi.org/10.1183/13993003.01200-2024}{\textit{Eur Respir J}. 2024;2401200.}

    \item Waxman A, Restrepo-Jaramillo R, Thenappan T, et al. Inhaled treprostinil in pulmonary hypertension due to interstitial lung disease. \href{https://doi.org/10.1056/NEJMoa2008470}{\textit{N Engl J Med}. 2021;384:325-334.}

    \item Nathan SD, Behr J, Collard HR, et al. Riociguat for idiopathic interstitial pneumonia-associated pulmonary hypertension (RISE-IIP): a randomized, placebo-controlled phase 2b study. \href{https://doi.org/10.1016/S2213-2600(19)30250-4}{\textit{Lancet Respir Med}. 2019;7(9):780-790.}

    \item Raghu G, Behr J, Brown KK, et al. Treatment of idiopathic pulmonary fibrosis with ambrisentan: a parallel, randomized trial. \href{https://doi.org/10.7326/0003-4819-158-9-201305070-00003}{\textit{Ann Intern Med}. 2013;158(9):641-649.}
\end{enumerate}

\end{document}
